% Copyright (c) 2026 Eugene Joseph M. Ragasa.
% SPDX-License-Identifier: Apache-2.0
% Author: Eugene Joseph M. Ragasa
% Formal proof completed: 2026-10-20
\documentclass[11pt]{article}
\usepackage[T1]{fontenc}\usepackage{amsmath,amssymb}\usepackage[margin=1in]{geometry}\usepackage[hidelinks]{hyperref}
\title{Proof Dossier: Idempotence of a Normalized Rank-One Projector}\author{Eugene Joseph M. Ragasa}\date{Formal proof completed: 2026-10-20}
\begin{document}\maketitle
\section{Identity and status}
ID \texttt{QM-DIRAC-THM-001}; status \texttt{proved}. Lean declaration:
\par\noindent{\scriptsize\path{ProjectKoios.QuantumMechanics.pureStateProjector_idempotent}}.
\section{Exact statement}
For a unit vector $\psi$ in a complex inner-product space and
$P_\psi=\lvert\psi\rangle\langle\psi\rvert$,
\[P_\psi\circ P_\psi=P_\psi.\]
The outer product is the bounded rank-one operator $v\mapsto\langle\psi,v\rangle\psi$.
\section{Source and assumptions}
This is standard source-neutral Hilbert-space mathematics. Exact textbook attachment is deferred.
The theorem concerns an ordinary vector, not a generalized ket or an unbounded operator.
\section{Proof schematic and proof}
For every $v$,
$P_\psi(P_\psi v)=\langle\psi,v\rangle\langle\psi,\psi\rangle\psi$.
Unit norm gives $\langle\psi,\psi\rangle=1$, leaving $P_\psi v$.
\section{Lean realization and trusted base}
The proof specializes Mathlib's idempotence theorem for \path{rankOne} and rewrites bounded-map
multiplication as composition. Kernel-reported axioms are \path{propext},
\path{Classical.choice}, and \path{Quot.sound}; there are no admissions.
\section{Boundary and review}
The theorem alone does not establish self-adjointness, a spectral decomposition, physical state
preparation, or source fidelity. AI-assisted and kernel-checked; no independent expert review is
recorded.
\end{document}
